\section{Full proof appendices}
\label{app:proofs}

This appendix records the complete finite-dimensional arguments used by the
claim registry. No statement below depends on an optimizer, a conjecture, or
a numerical plot.

\subsection{Magnitude induction and exact root geometry}

For a subtree rooted at \(v\), let \(k_v\) be its number of internal
vertices and let \(L_v\) be the product of the reduced leaf norms below it.

\begin{lemma}[Subtree magnitudes]
\label{lem:magnitudes}
If every node law has multilinear operator norm at most \(M\), then
\[
  \lVert F_v\rVert\le M^{k_v}L_v,
  \qquad
  \lVert R_v\rVert\le M^{k_v}L_v.
\]
\end{lemma}

\begin{proof}
At a leaf, \(F_v=R_v=Q_vz_v\), and \(Q_v\) is an isometry. Suppose the
claim holds for the children \(c_1,\ldots,c_a\). Multilinearity and the
definition of the operator norm give
\[
 \lVert F_v\rVert
 \le M\prod_i \lVert F_{c_i}\rVert
 \le M^{1+\sum_i k_{c_i}}\prod_iL_{c_i}
 =M^{k_v}L_v.
\]
The same argument applies to
\(R_v=P_v\mu_v(R_{c_1},\ldots,R_{c_a})\), because an orthogonal projector
is contractive. Induction proves both claims.
\end{proof}

\begin{proposition}[Root-error orthogonality]
\label{prop:rootgeometry}
For every valid typed tree,
\[
 (E_T^{\rm amb})^2=(E_T^P)^2+(E_T^N)^2,
 \qquad E_T^{\rm red}=E_T^P.
\]
\end{proposition}

\begin{proof}
At the root, write \(P=QQ^*\), \(F=F_T\), and
\(R=R_T\in\operatorname{ran}P\). Then
\[
 F-R=(PF-R)+(I-P)F.
\]
The first summand belongs to \(\operatorname{ran}P\) and the second to its
orthogonal complement. Pythagoras gives the first identity. Moreover,
\[
 Q^*F-Q^*R=Q^*(PF-R).
\]
The restriction of \(Q^*\) to \(\operatorname{ran}P\) is an isometry,
which proves the second identity.
\end{proof}

\subsection{Exact subset expansion}

\begin{theorem}[Local child-error expansion]
\label{thm:subset-full}
Let \(v\) have arity \(a\), children \(c_i\), recursively projected values
\(R_i\), and errors \(D_i=F_i-R_i\). For
\[
 y_i^S=\begin{cases}D_i,&i\in S,\\R_i,&i\notin S,\end{cases}
 \qquad
 r_v=(I-P_v)\mu_v(P_{c_1}\,\cdot,\ldots,P_{c_a}\,\cdot),
\]
one has the exact identity
\[
 D_v=r_v(R_1,\ldots,R_a)
 +\sum_{\varnothing\ne S\subseteq[a]}
 \mu_v(y_1^S,\ldots,y_a^S).
\]
\end{theorem}

\begin{proof}
Since \(F_i=R_i+D_i\), repeated multilinearity yields
\[
 \mu_v(F_1,\ldots,F_a)
 =\sum_{S\subseteq[a]}\mu_v(y_1^S,\ldots,y_a^S).
\]
The term indexed by the empty set is \(\mu_v(R_1,\ldots,R_a)\), while
\(R_v=P_v\mu_v(R_1,\ldots,R_a)\). Subtracting \(R_v\) separates the
empty-set term into its normal component \(r_v(R_1,\ldots,R_a)\); all
nonempty terms remain unchanged. No inequality has been used.
\end{proof}

Applying \(P_v\) to Theorem~\ref{thm:subset-full} removes the local
residual exactly. Applying \(I-P_v\) retains it. Decomposing each \(D_i\)
into its projected and normal parts produces the three-state alphabet
\(\{R,D^P,D^N\}\) used by the mixed-mask dynamic program.

\subsection{Universal homogeneous constants}

\begin{theorem}[Universal \(k\) and \(k-1\) estimates]
\label{thm:homogeneous-full}
Suppose \(\lVert\mu_v\rVert_{\rm op}\le M\) and
\(\lVert r_v\rVert_{\rm op}\le\rho\) at every internal vertex. If \(T\)
has \(k\) internal vertices and reduced leaf-norm product \(L_T\), then
\[
 E_T^{\rm amb}\le k\rho M^{k-1}L_T,
 \quad
 E_T^P=E_T^{\rm red}\le(k-1)\rho M^{k-1}L_T,
 \quad
 E_T^N\le k\rho M^{k-1}L_T.
\]
\end{theorem}

\begin{proof}
The claims are trivial for a leaf. At an internal vertex \(v\),
multilinear telescoping in any fixed slot order gives
\begin{align*}
 &\lVert\mu_v(F_1,\ldots,F_a)-\mu_v(R_1,\ldots,R_a)\rVert\\
 &\quad\le
 M\sum_{j=1}^a \lVert F_j-R_j\rVert
 \prod_{i<j}\lVert R_i\rVert
 \prod_{i>j}\lVert F_i\rVert.
\end{align*}
Assume the ambient estimate for every child.
Lemma~\ref{lem:magnitudes} shows that the \(j\)-th summand is at most
\(k_{c_j}\rho M^{k_v-1}L_v\). Because
\(\sum_j k_{c_j}=k_v-1\), the total propagated-child contribution is at
most \((k_v-1)\rho M^{k_v-1}L_v\). The local residual is bounded by
\[
 \rho\prod_j\lVert R_j\rVert
 \le \rho M^{k_v-1}L_v.
\]
Their sum proves the ambient estimate. For the projected error,
\[
 P_vD_v=P_v\bigl(\mu_v(F_1,\ldots,F_a)
                  -\mu_v(R_1,\ldots,R_a)\bigr),
\]
so the root local residual is absent and only the coefficient \(k_v-1\)
remains. Proposition~\ref{prop:rootgeometry} identifies the reduced error
with this projected error. For the normal error, the local residual is
retained, and the ambient argument gives coefficient \(k_v\). Induction
closes the proof.
\end{proof}

\subsection{Optimal telescoping order}

Let \(e_i,r_i,f_i,G_i\ge0\), set \(w_i=G_ie_i\), and for a permutation
\(\pi\) define
\[
 C(\pi)=\sum_t w_{\pi_t}
 \prod_{s<t}r_{\pi_s}\prod_{s>t}f_{\pi_s}.
\]

\begin{theorem}[Pair-exchange ordering]
\label{thm:ordering-full}
Put \(d_i=f_i-r_i\). A global minimizer of \(C\) is obtained by placing
the slots with \(d_i>0\) first, then those with \(d_i=0\), then those with
\(d_i<0\), and sorting each nonzero sign class by nondecreasing
\(w_i/d_i\). Slots with \(w_i=d_i=0\) are indifferent.
\end{theorem}

\begin{proof}
Consider adjacent slots \(i,j\). Prefix and suffix products are common to
the two possible orders. Thus \(i\) before \(j\) is no worse precisely when
\[
 w_if_j+r_iw_j\le w_jf_i+r_jw_i
 \quad\Longleftrightarrow\quad
 w_id_j\le w_jd_i.
\]
If both \(d_i,d_j\) are nonzero and have the same sign, division by the
positive product \(d_id_j\) gives \(w_i/d_i\le w_j/d_j\). If \(d_i>0\)
and \(d_j<0\), the exchange inequality holds automatically, so every
positive-denominator slot precedes every negative one. The cases with one
zero denominator place the zero class between them. Repeatedly exchanging
inverted adjacent pairs never increases \(C\) and terminates at the stated
order. Every permutation can be transformed in this way, proving global
optimality within the declared telescoping family.
\end{proof}

\subsection{Mixed-mask and path-sum certificates}

\begin{theorem}[Sound dynamic program]
\label{thm:dp-full}
Assume certified bounds for every full, projected-output, and normal-output
mixed block of each node law. The recursion retaining
\(B_v^F,B_v^R,B_v^A,B_v^P,B_v^N\) and all non-all-\(R\) state vectors in
\(\{R,D^P,D^N\}^{a_v}\) returns valid upper bounds. Its work is
\[
 O\!\left(|T|3^{a_{\max}}+|T|a_{\max}\log a_{\max}\right).
\]
\end{theorem}

\begin{proof}
At a leaf all five bounds are exact. At an internal node, expand by
Theorem~\ref{thm:subset-full} and then split every child error orthogonally
into its \(D^P\) and \(D^N\) parts. Every resulting term has one of three
states in each slot. Multiplying the corresponding certified mixed-block
norm by the child magnitude bounds and summing all state vectors except
all \(R\) is a direct application of the operator-norm inequality. The
local residual contributes
\(\lambda_v=\rho_v\prod_iB_{c_i}^R\) only to full and normal output. Taking
the minimum of independently valid upper bounds preserves validity.

There are \(3^{a_v}-1\) nontrivial state vectors at node \(v\), while
Theorem~\ref{thm:ordering-full} requires sorting \(a_v\) slots. Summing
these costs over the internal vertices proves the complexity estimate.
\end{proof}

\begin{corollary}[Residual-source path sum]
\label{cor:path-full}
Fix a valid telescoping order at every ancestor. If \(h_{a,j}\) is the
certified gain for propagating an error from child slot \(j\) through
ancestor \(a\), then
\[
 B_{T,\rm path}^{\rm amb}
 =\sum_{v\in\operatorname{Int}T}\lambda_v
   \prod_{(a,j)\in\operatorname{path}(v,\operatorname{root})}h_{a,j}
\]
is valid. For projected-root error the root source is omitted and the last
gain is a projected-output gain.
\end{corollary}

\begin{proof}
The scalar nodewise recurrence is affine in each local source
\(\lambda_v\). Repeated substitution sends that source through exactly the
ancestor edges on its root path and multiplies their gains. Summing over
all possible sources gives the formula. Projection annihilates the root
normal source, which proves the projected variant.
\end{proof}

\subsection{Signed forests and representation perturbations}

\begin{corollary}[Signed-forest triangle certificate]
If \(\mathcal F=\sum_\alpha c_\alpha T_\alpha\) is a finite compatible
signed forest, then the projected discrepancy is bounded by the sum of the
per-tree projected certificates weighted by \(\lvert c_\alpha\rvert\).
In particular, the two-term, two-node ternary associator has coefficient
at most two.
\end{corollary}

\begin{proof}
Apply Theorem~\ref{thm:homogeneous-full} to each tree and then the triangle
inequality. Each two-node tree contributes \(k-1=1\), hence two signed
terms contribute at most two. Syntactically identical trees may first be
combined exactly.
\end{proof}

\begin{proposition}[Approximation plus projection]
\label{prop:cp-full}
Suppose \(\lVert\mu_v\rVert\le M\),
\(\lVert\mu_v-\widehat\mu_v\rVert\le\delta\), and
\(\lVert\widehat\mu_v\rVert\le M+\delta\). For a \(k\)-node homogeneous
tree,
\[
 E_{\rm repr}\le k\delta(M+\delta)^{k-1}L,
\]
and a closure coefficient \(c\in\{k,k-1\}\) may be split as
\[
 c\rho(M+\delta)^{k-1}L
 =c\rho M^{k-1}L
 +c\rho\bigl[(M+\delta)^{k-1}-M^{k-1}\bigr]L.
\]
\end{proposition}

\begin{proof}
Insert the full ambient \(\widehat\mu\)-tree between the exact ambient tree
and the recursively projected approximate tree. Multilinear telescoping
over the \(k\) node-law replacements bounds each replacement by
\(\delta(M+\delta)^{k-1}L\), giving the representation term. Apply
Theorem~\ref{thm:homogeneous-full} to the approximate law and add and
subtract \(c\rho M^{k-1}L\); the displayed closure and interaction terms
are an algebraic partition of that valid upper bound. No equality with the
observed total error is asserted.
\end{proof}
